% !TEX root = ../monograph.tex
\chapter[Local matching and the one-loop coefficient]{Local matching and the one-loop coefficient}\label{ch07:matching}

A renormalization step has two distinct normalization tasks. Its
classical action at the critical configuration must retain the prescribed quadratic form
and agree with a gauge-invariant local action through degree four.
Its Gaussian determinant changes the coefficient of that quadratic
form. The first task is a local Taylor construction; the second requires
a comparison of determinants uniform in the block ratio. This chapter
separates the two, so that a perturbative coefficient does not stand in
for an estimate on the block integral.

Fix one step $m+1\to m$, an odd block ratio $L=L_m$, and the Lie
algebra metric of Chapter~1. Momentum is measured in lattice units at
the indicated scale. The adjoint Casimir $C_G$ is defined by
$\tr_{\gLie}(\operatorname{ad}X\operatorname{ad}Y)
=-C_G\langle X,Y\rangle$; set
\[
 \kappa_G=\frac{11C_G}{24\pi^2}.
\]
The group is compact with simple nonabelian Lie algebra. Constants may
depend on fixed stage bounds, but constants described as uniform in $L$
are chosen before $L$. All local Taylor expansions below concern the
analytic action near the flat field, where the indicators of the regularity condition are
constant. They do not differentiate at the boundary of the small-field region.

\section{Infinite-lattice Taylor kernels and Ward identities}\label{ch07:kernels}

Write $C_n(E)=D^nE(0)/n!$ for the normalized symmetric Taylor
coefficient of an action $E$. A translation-invariant coefficient is
specified by its values on plane waves. After division by the number
of lattice sites, its symbol is a tensor
$f(p_1,\ldots,p_{n-1})$, with $p_0=-\sum_{i>0}p_i$.
The tensor has one spacetime index and one argument in $\gLie$ for each field.
The symbol is periodic with period $2\pi$ in each independent momentum.
An \emph{infinite-lattice symbol} means one function chosen on the infinite lattice whose
restriction to each finite momentum grid describes the corresponding
torus coefficient. This requirement is stronger than a separate smooth
extension of each finite grid.

For $\sigma>0$, write $|f|_\sigma\le A$ when every component is
holomorphic and bounded by $A$ on the strip
$|\operatorname{Im}p_i^a|<\sigma$. There are only finitely many direction
and Lie algebra components in the basis of Section~\ref{ch01:gauge} at fixed $n$. Reducing the strip gives an
exponentially weighted sum of Fourier coefficients, with a constant
depending on $n,\sigma,G$. Thus strip bounds for infinite-lattice symbols imply locality
uniformly in the period. We say a symbol has a zero of order at least $r$
if all its Taylor coefficients of total degree below $r$ vanish at zero
external momenta.

Let $q_a(p)=e^{ip_a}-1$ and $d_1(p)v=q(p)\wedge v$ be the lattice
curl. Suppose an analytic gauge-invariant difference $E$ has vanishing
Taylor coefficients below degree $n$. Its infinitesimal gauge identity
at that degree contains only the linear gauge variation $d_0\omega$:
the nonlinear part multiplies lower coefficients, which are zero.
Consequently its first nonzero symbol satisfies
\begin{equation}\label{ch07:ward}
 \delta_i f=0\quad(0\le i<n),\qquad
 \delta_i=\iota_{q(p_i)}.
\end{equation}
Contraction is algebraic, without complex conjugation; it acts only on
the direction index of the indicated \emph{leg} (field argument of the symbol). Reality means
$\overline{f(-p)}=f(p)$ on real momenta. Constant gauge transformations
make the tensor invariant under the diagonal adjoint action. A Ward identity verified on
all finite momentum grids passes to the continuous infinite-lattice symbol by
density. This last step requires an infinite-lattice kernel.

\section{Division by the lattice curl}\label{ch07:division}

The algebraic Poincar\'e lemma explains why gauge-invariant leading
coefficients are built from curvature; see
\cite[Section~4]{BarnichBrandtHenneaux2000}. Its local continuum form does
not give the periodic strip bounds needed here. We give the quantitative
adaptation, including the constraint on the sum of momenta.

\begin{lemma}[Quantitative curvature representation]\label{ch07:curl}
Let $f$ have $n$ one-form legs, satisfy \eqref{ch07:ward}, and obey
$|f|_\sigma\le A$. For $n\ge3$ there are $\eta>0,K<\infty$,
depending only on $n,\sigma$, and a periodic holomorphic tensor $g$
with two-form legs such that
\[
 f=\delta_0\cdots\delta_{n-1}g,\qquad |g|_\eta\le KA.
\]
For $n=2$ the same assertion holds if $f$ has order $r+2$, $r\ge0$;
then $g$ can be chosen with order $r$. Reality and adjoint-invariant
coefficients can be imposed with a fixed additional norm factor.
\end{lemma}

\begin{proof}
The elementary division is
\[
 D_af(p)=\frac{f(p)-f(p|_{p_a=0})}{e^{ip_a}-1}.
\]
The expression $p|_{p_a=0}$ means that only the $a$th coordinate is set
to zero. Periodicity removes the apparent singularities at $p_a\in2\pi\ZZ$.
For the Fourier monomial $e^{iz\cdot p}$ the quotient is a finite
geometric sum with $|z_a|$ terms and no frequency larger
than $|z_a|$. A smaller strip absorbs this polynomial factor into the
exponential Fourier decay. Hence $D_a$ is bounded from a fixed strip to
a smaller one and loses at most one Taylor order.

For a cycle $\delta f=0$ in one exterior factor, write
$f=e_a\wedge u+v$, with $u,v$ containing no $e_a$. Write $\delta'$ for contraction with $q$ with its $a$th component omitted.
Then $\delta'v=-q_au$, and direct substitution gives
\[
 f=\delta(e_a\wedge D_av)+v|_{p_a=0}.
\]
Repeat over the four axes. At positive exterior degree the last
remainder has no possible component and is zero. At degree zero the
last remainder is $f(0)$, so that case is solvable precisely when
$f(0)=0$. The construction preserves auxiliary constraints independent
of the momentum being divided. It raises the exterior degree by one,
loses at most one Taylor order, and has a bounded smaller-strip norm.

Only the dependent momentum $p_0$ needs extra care. The following
induction gives it. With $t$ independent momenta and exterior degrees
$(d,1,\ldots,1)$, a cycle in every factor has a zeroth-factor primitive
preserving the other cycles whenever $d+2t>4$.
For $t=0$ its degree exceeds four and it is zero. For $t>0$, replace
$p_1$ by $p_0$ as independent variable and apply the preceding division.
This invertible integral shear only reduces the strip by a fixed factor.
It gives $\delta_0X=f$, preserving the cycles in factors $2,\ldots,t$.
Apply another zeroth-factor division to $\delta_1X$, obtaining
$\delta_0Y=\delta_1X$. At $p_1=0$ the tensor $Y_0$ is a cycle with
zeroth degree $d+2$ and $t-1$ independent momenta. Induction gives
$\delta_0Z=Y_0$. Extend $Z$ independently of $p_1$ and set
$Y'=Y-\delta_0Z$. It vanishes at $p_1=0$, so scalar division in factor
one gives $\delta_1D=Y'$. The corrected primitive
$X-\delta_0D$ has contraction $f$ in factor zero and zero in factor one.
All other cycles were retained. The invariant quantity $d+2t$ makes
this a terminating induction.

Initially $d=1,t=n-1$, so the induction applies for $n\ge3$.
After the zeroth-factor primitive, divide the independent legs successively.
For $n=2$, start instead with an order-$r+2$ tensor. The first primitive has
order $r+1$; contraction in factor one restores one order, and the next
primitive $Y$ has order $r+1$. Thus $Y(0)=0$ without induction. Its scalar
primitive has order $r$, the corrected primitive has order $r+1$, and the
last division leaves order $r$. All operations are bounded on a common
smaller strip and their number is fixed.

Finally replace $g$ by $(g(p)+\overline{g(-p)})/2$ and project each
coefficient tensor on $\gLie$ onto the invariant real tensors. Contraction acts
on direction indices and commutes with this finite-dimensional
projection. These operations preserve the reconstruction and all
retained orders, with a fixed norm factor.
\end{proof}

Choose coordinate contours from the origin to the Fourier displacements
$z_1,\ldots,z_{n-1}$. The transported curvature along such a
contour is
\begin{equation}\label{ch07:slot}
 F_{\gamma,ab}(U,x)=\Ad(U_{x+\gamma})
 \operatorname{proj}_{\gLie}\frac{U_{x+z;ab}-U_{x+z;ab}^{-1}}2,
 \qquad \gamma:0\longrightarrow z.
\end{equation}
All factors transform at the same base point $x$, so an invariant
multilinear form on $\gLie$ is gauge invariant. Its length is the
sum of the contour lengths plus four per plaquette. The linearization of
each factor is the corresponding curl. The Fourier coefficients of
$g$ therefore produce an exact leading Taylor coefficient $f$.
Symmetrizing the legs implements the convention $C_n=D^n/n!$;
it introduces no extra factorial.

If $a$ denotes such a multilinear term and $|a|$ the sum of absolute
components in the basis of Section~\ref{ch01:gauge}, the result is
\begin{equation}\label{ch07:path-norm}
 \sum_a |a|e^{w\operatorname{length}(a)}\le K_w A
 \qquad(0<w<\eta).
\end{equation}
Averaging over signed coordinate symmetries preserves matching and this
bound. Absolute convergence permits periodization on every torus.
For the two-leg case with $r=2$, the curvature coefficients $c(z)$ obey
\begin{equation}\label{ch07:moments}
 \sum_zc(z)=0,\qquad \sum_z z_ac(z)=0\quad(1\le a\le4).
\end{equation}
These are the value and first derivatives of its symbol at zero.
They hold separately for each orientation and each pair of Lie algebra arguments.

\section{The marginal quadratic coefficient}\label{ch07:quadratic}

A quadratic Taylor symbol includes the factor $1/2!$. For its transverse
entry, define
\begin{equation}\label{ch07:alpha}
 \alpha(f)=2\operatorname{Re}[t^2]f_{22}^{uu}(te_1),
\end{equation}
where $[t^2]$ extracts the coefficient of $t^2$ at $t=0$, and the
superscript $uu$ means that the same unit vector $u\in\gLie$ is inserted in
both arguments. The quadratic homogeneous term of the kinetic symbol $b_H$
of $J$, evaluated on $v,w\in\gLie$, is
\[
 (b_H)_{ab}^{vw,[2]}(p)=\tfrac12
 (|p|^2\delta_{ab}-p_ap_b)\langle v,w\rangle_{\gLie},
\]
so $\alpha(b_H)=1$. The factor two in \eqref{ch07:alpha} is therefore
fixed by the Taylor convention, not by a choice of coupling scheme.

\begin{lemma}[Removal of the kinetic coefficient]\label{ch07:unique}
An analytic quadratic symbol from an invariant real action, with
exchange symmetry $f_{ab}^{uv}(p)=f_{ba}^{vu}(-p)$ and signed-coordinate symmetry, has
$f-\alpha(f)b_H=O(|p|^4)$. Its coefficient $\alpha$ is independent
of the transverse axis pair and the unit vector $u$.
\end{lemma}

\begin{proof}
The Ward identity first annihilates the constant term. Its linear term is
also zero: inversion of all coordinates acts positively on the two
one-form indices and negatively on a linear momentum polynomial.
Edge translation phases have constant term one; their higher terms
multiply the already vanishing constant coefficient.

At quadratic order the symmetric diagonal forms on $\gLie$ are scalar
multiples of the invariant metric. A self-adjoint operator commuting
with the adjoint action has invariant eigenspaces, which are ideals;
simplicity leaves only one eigenvalue. Coordinate reflections allow
only $p_c^2$ in a diagonal direction entry and $p_ap_b$ in an
$a,b$ off-diagonal entry. Permutation symmetry identifies their
coefficients. The transverse Ward equations then force the tensor to
be a scalar multiple of $|p|^2I-pp^{\mathsf T}$.
One transverse evaluation determines that scalar. Subtracting the
kinetic symbol $b_H$ itself preserves all symmetries, not merely their leading
parts. Its residual cubic coefficient vanishes by inversion, since all
lower coefficients now vanish. Thus the residual has order four.
\end{proof}

Apply the two-leg curvature representation to this residual. The
resulting quadratic curvature terms have the zero moments
\eqref{ch07:moments}. They cannot change the coupling: their curvature
symbol is $O(|p|^2)$ and their link symbol is $O(|p|^4)$.
The coefficient of curvature square is carried entirely by $\beta$.
This statement refers to Taylor coefficients at the flat field; it does
not identify a nonlinear effective action with its quadratic truncation.

\section{Haar measure and the gauge-fixing term}\label{ch07:completion}

The Gaussian contribution at a critical configuration $U_v$ of the fine field is
\begin{equation}\label{ch07:thermal}
 T(v)=\mathcal L_{m+1}(U_v)+\tfrac12\log\det Q(v)-\log J_{\rm Haar}(v).
\end{equation}
Here $Q$ is the fluctuation Hessian without its factor $\beta$,
$J_{\rm Haar}$ is the density of Haar measure in the chosen chart, and
$\mathcal L_{m+1}$ is the $\beta$-independent quadratic term of the path action, called the \emph{inherited loop action}.
Field-independent constants may be subtracted throughout. This does
not alter any Taylor coefficient of positive degree.

The following finite-dimensional identity accounts for the ghost
factor. Let $\mathsf J$ be the joint Hessian in left tangent coordinates,
$\mathcal R$ the infinitesimal gauge action, and $\mathcal I$ a
slice embedding. Assume
\[
 \mathsf J\mathcal R=0,\quad |\det(\mathcal R,\mathcal I)|=1,
 \quad Q=(Ds)^*\mathcal I^*\mathsf J\mathcal I\,Ds,
 \quad J_{\rm Haar}=|\det Ds|.
\]
These are identities at the critical configuration. Given a gauge
condition $F$ with invertible $G=F\mathcal R$, put
$C=\mathsf J+F^*F$. Block triangular factorization in the basis
$(\mathcal R,\mathcal I)$ gives
\[
 \det C=(\det G)^2\det(\mathcal I^*\mathsf J\mathcal I),
 \qquad
 \tfrac12\log\det Q-\log J_{\rm Haar}
 =\tfrac12\log\det C-\log|\det G|.
\]
The fluctuation form is positive on the slice. With a positive regulated
ghost operator the absolute value can be dropped. This is the
finite-dimensional Faddeev--Popov determinant identity; the moving
Haar Jacobian is part of it. Omitting that Jacobian would change the
quantity being compared. Background-field perturbation theory uses
exactly this vector-minus-ghost combination~\cite{Abbott1982}.

For the lattice comparison choose $\mu=L^{-1}$ and
\[
 P_\mu=\mu^2(\Delta_U+\mu^2I)^{-1},\qquad
 m_\mu=\mu^2P_\mu^9,\qquad \Delta_U=d_U^*d_U.
\]
If $X\mathcal R=I$ is the derivative of the local gauge transformation, take
$F=d_U^*P_f+m_\mu X$; then $G=\Delta_U+m_\mu$.
The regulated fine operator is
$A_\mu=A+d_Ud_U^*+M_\mu$, where $A$ is the fine kinetic Hessian
and $M_\mu$ applies $m_\mu$ to each one-form component.
Its determinant response is
$\Pi_\mu=\frac12\log\det A_\mu-\log\det G$.
The same ghost appears in both expressions and cancels in their
difference.

\section{The regulated lattice determinant}\label{ch07:one-loop}

Perturbative matching for the Wilson lattice action is treated in
L\"uscher--Weisz~\cite{LuscherWeisz1995}; our comparison additionally
controls the selected non-Wilson path action and its averaging operation.
The continuum coefficient is the classical calculation of
Gross--Wilczek and Politzer~\cite{GrossWilczek1973,Politzer1973}.
We use the background-field formulation~\cite{Abbott1982}; the following
normalization check and integrable error criterion are the parts needed
for the lattice application. Recall that $\alpha$ includes the factor
two converting a normalized quadratic Taylor coefficient into the
coefficient of $\frac12\langle d_1v,d_1v\rangle$.

Let $k$ denote the loop momentum and $E_{ab}$ the matrix unit on
spacetime indices. For external momentum $te_1$, polarization $e_2$,
and unit vector $u\in\gLie$, the leading ghost vertex is
$-2ik_2\operatorname{ad}u$. The vector
vertex has four copies of this orbital term and the spin term
$2it(E_{21}-E_{12})\operatorname{ad}u$. The orbital--spin trace is zero.
With vector and ghost bubble coefficients $-1/4$ and $+1/2$, respectively,
and $\tr(\operatorname{ad}u)^2=-C_G$, the coefficient extracted by
\eqref{ch07:alpha} has homogeneous integrand
\begin{equation}\label{ch07:pole}
 h(k)=4C_G\left(\frac{k_2^2}{|k|^6}
       -\frac{4k_1^2k_2^2}{|k|^8}-\frac1{|k|^4}\right).
\end{equation}
This is the standard background-field bubble calculation, with its
Lie algebra normalization expressed in our metric. It is not a claim that
the full lattice integrand equals $h$.

\begin{proposition}[The lattice integrand remainder]\label{ch07:integrand-input}
Let the incoming analytic path action have a fixed positive exponential
coefficient bound, a uniformly positive curvature metric, and kinetic
normalization $H(0)=I$. Form its left-coordinate Hessian (the Hessian in left increments of the proof below), add
the gauge-fixing term $d_Ud_U^*$ to it, and use the covariant regulator $m_\mu$ defined above.
The regulated quadratic coefficient has an integrand $I_\mu$ satisfying
\begin{align*}
 \alpha(\Pi_\mu)&=(2\pi)^{-4}\int_{[-\pi,\pi]^4}
                                    \operatorname{Re}I_\mu(k)\,dk,\\
 |I_\mu(k)|&\le A(|k|+\mu)^{-4},\\
 |I_\mu(k)-h(k)|&\le D(|k|^{-3}+\mu^2|k|^{-6})
                  \quad(\mu\le|k|\le1).
\end{align*}
The constants depend on the fixed path bounds, on the lower bound for the metric and on the group,
but not on $\mu\in(0,1]$ or the torus period. The finite-torus Hessian
is the corresponding momentum mesh average, with an exponentially small
per-site error relative to the integral at each fixed $\mu$.
\end{proposition}

\begin{proof}
We specify the derivatives used in the comparison. In one vector or
ghost channel let $a(k)$ be the zero-background inverse propagator,
$b(p;k)$ its first derivative in the background field, for short its first \emph{background derivative}, and $c(p,p';k)$ its second
background derivative. These are exact derivatives of the same operator; the
internal momentum of $b(p;k)$ changes from $k$ to $k+p$.
Their contribution to the logarithmic Hessian is
\begin{equation}\label{ch07:response-integrand}
 \tr\bigl(a(k)^{-1}c(-p,p;k)
       -a(k)^{-1}b(-p;k+p)a(k+p)^{-1}b(p;k)\bigr).
\end{equation}
Take half the vector expression minus the ghost expression and apply
the normalized external coefficient functional. This defines $I_\mu$.
Finite Fourier inversion gives the identical formula with the internal
integral replaced by the mesh mean. Thus the integrand is associated
with the fine action, rather than with a separately chosen
continuum model.

We first check its local Taylor data. A transported plaquette curvature
has, under the formal scaling $A(x)=\epsilon a(\epsilon x)$,
leading term $\epsilon^2(da+[a,a])$ in the corresponding orientation.
Its remaining terms have scaling dimension at least three; a Taylor derivative or
a factor of $A$ counts as one unit. This follows by expanding the four
edge exponentials and their parallel transports. At scaling dimension four the row squares
and direct term therefore have coefficient on $\gLie$-valued one-forms
$c_{\rm dir}I+\widehat\rho(0)^*\widehat\rho(0)=I$.
Cubic and quartic curvature terms have scaling dimensions at least six and eight.
The leading local density is consequently the Wilson curvature-square
density. All differences from it have scaling dimension at least five.
Exponential path summability justifies these expansions term by term
and bounds their Taylor remainders: each derivative contributes a polynomial
in path length, absorbed by a smaller \emph{width}, that is, a smaller exponential decay rate of the path coefficients.

The Hessian in left increments is obtained from
$S(\log(e^{x_+}e^{x_-}))$ by two $x_+$ derivatives. The
Baker--Campbell--Hausdorff correction starts with two fields and hence
raises the scaling dimension relative to its linear term. Accordingly, after two
fluctuation and $n$ background derivatives, the difference from the Wilson
response has momentum order at least $3-n$, for $0\le n\le2$.
This is a statement about the combined momenta of all legs. It gives
uniform compact derivative bounds as well as the indicated orders.

In particular the gauge-fixed vector inverse propagator is
\[
 a_{\rm loc}(k)=d_1(k)^*H(k)d_1(k)+d_0(k)d_0(k)^*
              =|k|^2I+O(|k|^3).
\]
The ghost inverse propagator is $\omega(k)I$, with
$\omega(k)=\sum_a|q_a(k)|^2=|k|^2+O(|k|^4)$.
For a mark of direction $a$ and Lie algebra direction $u$, expanding the Wilson cubic
vertex and adding the gauge-fixing vertex gives the linear part of the first vector
background derivative
\[
 b_{\rm lin}(p;k)_{ij}
 =\bigl[-i(2k_a+p_a)\delta_{ij}
              +2ip_j\delta_{ia}-2ip_i\delta_{ja}\bigr]
                         \operatorname{ad}u.
\]
The linear part of the first ghost background derivative is $-i(2k_a+p_a)\operatorname{ad}u$.
These expressions can also be read from the continuum quadratic
operator in background gauge. The lattice remainder in each is
$O((|k|+|p|)^2)$, by the comparison of scaling dimensions just proved. The second
background derivative and its first two external momentum derivatives are uniformly
bounded on the real momentum cube. No cancellation of a seagull term
has been assumed.

For completeness we give the bounds that propagate this information
through the regulator and the inverses. At a fixed internal momentum
$k$ put $T=|k|+\mu$, and for a matrix path $M(t)$ write
\[
 \|M\|_{T,2}=\max_{0\le j\le2}T^j\|M^{(j)}(0)\|.
\]
The paths have displaced momenta $k+te_1$, with velocities of bounded
size. The product rule bounds products by a fixed multiple of the
product of these norms. The two inverse derivative identities give
$\|M^{-1}\|_{T,2}\le C/T^2$ whenever
$\|M\|_{T,2}\le CT^2$ and $\|M(0)^{-1}\|\le C/T^2$.

The lower bound on the inverse propagators at real momenta holds globally on the centered cube.
Indeed $\omega(k)\ge c|k|^2$, the exterior identity controls the sum of
transverse and longitudinal parts, and
\[
 x+\mu^2\left(\frac{\mu^2}{x+\mu^2}\right)^9
                    \ge2^{-10}(x+\mu^2)\quad(x\ge0).
\]
Thus the inverse propagators of both channels are bounded below by $cT^2I$.
The massive ghost operator $\Delta_U+\mu^2I$ has bounds for its inverse propagator and its
first and second background derivatives, of respective sizes
$T^2,T,1$ in the scaled derivative norms. Inverting gives sizes
$T^{-2},T^{-3},T^{-4}$: a first background derivative is $-R bR$, and a second
is the sum of $-RcR$ and the two ordered $RbRbR$ terms.
Multiplication by $\mu^2$, taking nine products, and multiplying once
more by $\mu^2$ gives regulator sizes
\begin{equation}\label{ch07:regulator-budgets}
 \mu^2,\qquad \mu^2/T,\qquad \mu^2/T^2.
\end{equation}
This uses $\mu^2/T^2\le1$. Product differentiation through order two
preserves the same sizes in $\|\cdot\|_{T,2}$; the number nine is fixed.
The four constant maps $J_a$ in the vector regulator $M_\mu=\sum_aJ_a^*m_\mu J_a$, where $J_a$ takes a one-form to its component in direction $a$, do not
change these powers.

The first term of \eqref{ch07:response-integrand}, differentiated twice
in $t$, is bounded by $CT^{-2}+C\mu^2T^{-6}$. Its unshifted inverse
is independent of $t$; the local second background derivative has a unit-scale
bound, while its regulator part uses \eqref{ch07:regulator-budgets}.
The second term has two inverses of size $T^{-2}$ and two first background derivatives
of size $T$, with two derivatives contributing $T^{-2}$. Its bound is
$CT^{-4}$. Since $T$ stays bounded above on the cube, these estimates
give the first asserted radial bound.

For the comparison use $T=|k|$ and restrict to $\mu\le T\le1$.
Let $\delta=T+\mu^2/T^2$. The regulated lattice inverse propagator minus the free
quadratic one has scaled norm at most $C\delta T^2$.
The resolvent identity therefore bounds the inverse difference by
$C\delta T^{-2}$. Likewise the first background derivative minus its displayed
linear part is bounded by $C\delta T$: the local remainder is
$O(T^2)$ and the regulator contribution is $O(\mu^2/T)$.
Subtract the free bubble by telescoping its four factors. Each term
contains one difference, so its second external derivative is bounded
by $C\delta T^{-4}=C(T^{-3}+\mu^2T^{-6})$.
The seagull estimate above is bounded by the same expression on this
annulus. The free vector-minus-ghost bubble is exactly
\eqref{ch07:pole}, proving the second radial inequality.

At fixed $\mu>0$ the inverse propagators are invertible on a smaller complex
strip by their positive real lower bound and analyticity. All three response
layers come from the same infinite-lattice path series. The integrand
thus has a common holomorphic extension for that fixed $\mu$.
For a periodic analytic function, the difference between its mesh mean
and its integral is the sum of nonzero Fourier coefficients whose
indices are multiples of the mesh side. Exponential coefficient decay
bounds this error by $C_\mu e^{-c_\mu n}$.
The constants in this last mesh estimate may depend on $\mu$;
the radial constants proved above do not. This proves the final assertion
and justifies the coefficient integral.
\end{proof}

\section{Uniform comparison of determinants}\label{ch07:comparison}

Two estimates must be distinguished. An absolute determinant Taylor
bound may grow like the number $L^4$ of fine sites per block. The
\emph{difference} between the blocked and regulated determinant must
have a bound independent of $L$ in order to identify the logarithmic
coefficient. The relevant mechanism is finite rank together with a
trace-class regulator.

Here is the structure behind the trace-norm requirement. Writing $H$
for the derivative of the averaging operation, $B$ for the Hessian of its penalty term, $Z$ for the
coarse-frame projection, and $K=m_\mu X$, cancellation of the fine
kinetic Hessians leaves
\begin{align}\label{ch07:theta}
 \Theta={}&H^*BH-Z^*Z+P_f^*d_UK+K^*d_U^*P_f+K^*K\\
 &\hspace{12mm}-P_f^*M_\mu P_f.\notag
\end{align}
The first two terms factor through spaces of dimensions at most
$4\dim\gLie$ and $\dim\gLie$ per coarse momentum. If $W$ maps into such an intermediate space of dimension $r$ and $V$
maps out of it, $VW$ has at most $r$ nonzero singular values. Thus its
trace norm satisfies $\|VW\|_1\le r\|V\|\|W\|$, uniformly in $L$.
Derivatives preserve the intermediate spaces.

For the regulator, at zero background the alias eigenvalues of $P_\mu$
are bounded by $C(1+|z|)^{-2}$ in a smaller complex strip. Consequently
\[
 \|P_\mu^3\|_1\le C\sum_{z\in\ZZ^4}(1+|z|)^{-6}<\infty.
\]
Given bounded first two background derivatives of $P_\mu$, distribute at most two
derivatives among the three factors in $(P_\mu^3)^3$. Every term retains
one undifferentiated factor $P_\mu^3$, which is trace class; the other factors have operator
bounds. Hence the first two background derivatives of $m_\mu$ have trace norm
$O(\mu^2)$. This explains the ninth power and avoids an erroneous
factor $L^4$. To apply this trace estimate, we establish the comparisons of weighted and Euclidean
norms, the background derivatives at the critical configuration and the scaled inverse bounds below.

\subsection{The critical configuration and the contact term of the averaging operation}
\label{ch07:contact}

We verify the assertion about the contact correction to the determinant, defined below, for the averaging operation used in the
block integral. Let $\zeta(x)$ be the representative of its critical configuration in centered
coordinates (in which the local gauge transformation is normalized in the mean, Section~\ref{ch07:weighted}), with $\zeta(0)=0$ and tangent $(Wx,r_x)$, where $r_x$ is its
coarse scalar component after recentering.
\mgubi{Issue 10: which $W$? Both $W_a$ (then $r_x=-RW_ax$) and $W_F$ are possible. Independence of the choice needs invariance of $f$ under fine gauge transformations; only constant ones are stated.}%
For each scalar determinant function $f(x,z)$, constant gauge transformations
act orthogonally on its matrices and leave $f$ invariant. Every invariant
linear functional on the simple Lie algebra vanishes, since it annihilates
brackets and $[\gLie,\gLie]=\gLie$. Applying this to each coordinate gives
$Df(0)=0$. Thus
\begin{equation}\label{ch07:tangent-chain}
 D^2[f(x,\zeta(x))](0)[u,v]
   =D^2f(0)[(u,Wu,r_u),(v,Wv,r_v)].
\end{equation}
Second derivatives of the critical-configuration map have disappeared for this
specific reason. Matrix Hessians themselves need not have vanishing first
derivative; the statement is about their invariant scalar determinants.

Use rooted coordinates (in which the local gauge transformation is normalized at the root of each block) for the following exact calculation. Let $w_x(v)$
be the logarithmic change in the averaging operation under a fluctuation left increment
$v$, with $w_x(0)=0$. Let $r(x)$ be the derivative of the penalty term
with respect to that increment of the averaging operation, and set
\[
 S_x(a)_{ij}=\langle a,D_v^2w_x(0)[e_i,e_j]\rangle,
 \qquad Z(x)=S_x(r(x)),\qquad A(x)=M_{\rm core}(x).
\]
The Hessian chain rule gives $M_{\rm act}=A+Z$, with $r(0)=Z(0)=0$; we call $Z$ the \emph{contact term} of the averaging operation, and the difference in \eqref{ch07:contact-exact} below the \emph{contact correction}.\mgubi{Issue 10: $M_{\rm act}$, $M_{\rm core}$ are not defined. From the Lean definitions: $M_{\rm act}$ is $C=J+F^*F$ at the critical configuration; in $M_{\rm core}$ the penalty Hessian is replaced by $H^*BH$ (probably $C_0$ up to the change of coordinates; not verified).}
The change of fluctuation coordinates from rooted to centered is triangular with determinant one,
so it does not change the determinant. Define $Y(x)$ by
$\langle a,Y(x)\rangle=\tr(A(x)^{-1}S_x(a))$.
The averaging operation is covariant, and the matrices on the fluctuation coordinates transform by the
same congruence. Cyclicity of trace therefore makes $Y$ a covariant
coarse vector. Constant gauge transformations imply $Y(0)=0$.

The linear block identity $QW=I-b^{-1}K_c$ fixes the derivative of the
penalty gradient:
$Dr(0)=-K_c=:J$. Expanding $\log\det(A+Z)-\log\det A$ through degree two
now gives
\begin{align}\label{ch07:contact-exact}
 D^2(\log\det M_{\rm act}-\log\det M_{\rm core})(0)[u,v]
 ={}&\langle Ju,DY(0)v\rangle+\langle Jv,DY(0)u\rangle\\
 &-\tr(A(0)^{-1}S_0(Ju)A(0)^{-1}S_0(Jv)).\notag
\end{align}
In particular the derivative of $A^{-1}$ is included in $DY$;
$Y(0)=0$ removes the term with $D^2r$.

There is an exact identity $DY(0)d_{0,c}=0$. Gauge the trivial
configuration by a coarse gauge transformation, then restore the chosen
block axial gauge. Gauge invariance preserves stationarity; uniqueness of the
local critical configuration identifies the restored configuration with the
critical configuration on this curve. Covariance and $Y(0)=0$ show that $Y$
vanishes there. Differentiation proves the identity. This uses the nonlinear
averaging operation, rather than an independently chosen linear averaging operator.

At fixed $L$, the finite-stencil averaging operation and the analytic inverse of
$A(0)$ give holomorphic infinite-lattice kernels for every term in
\eqref{ch07:contact-exact}. Normalized finite-torus traces are momentum
mesh averages. Their limit is the integral of these same kernels;
exponential Fourier decay gives uniform convergence also along varying
external mesh momenta. Write $M(p)$ for the mixed kernel of $DY$.
Its Ward identity along $p=te_\nu$ reads
$q_\nu(te_\nu)M_{j\nu}(te_\nu)=0$ as an identity on $\gLie$.
Hence every column of $M(0)$ is zero and $M(p)=O(|p|)$.
Since $J(p)=-K_c(p)=O(|p|^2)$, the first two terms of
\eqref{ch07:contact-exact} have order at least three, and the last has
order at least four. Therefore the contact correction has zero $\alpha$.
Only qualitative analyticity at fixed $L$ was used here. Its constants
need not be uniform in $L$, because this is an exact vanishing statement.

\subsection{Bounds for the centered matrices}\label{ch07:weighted}

The smooth representative of the linearization $W$ of the critical configuration is $W_a$ from \eqref{ch05:corrected-transport}; its representative in the block axial gauge is $W_F$ from \eqref{ch05:stationary-tangent}. As above, $\beta$ is the inverse coupling, and $b$ is the penalty parameter of the block kernel.

We establish the operator estimates for the selected path action and
averaging operation. The inputs are their positive normalized metric, fixed
exponential path bounds, and the identity \eqref{ch05:stationary-tangent} and smooth
representative of Section~\ref{ch05:transports}; all refer to the same
two consecutive scales. The coarse scalar part of the smooth representative
may be nonzero after recentering. Its derivatives below are merely
required to be bounded, not to vanish.

Recall that $\mu=L^{-1}$ is the inverse block ratio. Use the unitary
block Fourier transform $P(p)$ on positions within the block, and set
\[
 h_l(p)=(|q(\operatorname{Re}(p+2\pi l)/L)|^2+\mu^2)^{1/2},
 \qquad
 \|u\|_{p}^2=\sum_l h_l(p)^2|(P(p)^{-1}u)_l|^2.
\]
The joint space also has a coarse scalar component with its ordinary
Euclidean norm. Opposite momenta are paired in the dual norm.
On a fixed small coarse strip, $P$ and $P^{-1}$ have bounded operator
norms, $h_l\ge\mu$, and $|q(k_l)|\le5h_l$.
Thus a unit fluctuation direction has normalized $L^2$ norm over positions within the block
$O(\mu)$ and its first fine difference has normalized $L^2$ norm
$O(\mu^2)$. Background smooth bridge legs have the same powers in
supremum norm, by \eqref{ch05:actual-bridge}. Extending to the fixed
padded box of radius $(j+5)L$ around the patches of the averaging operation only changes constants: its number of residue
blocks is fixed at the given smoothing order, and its Bloch phases
are uniformly bounded.

\paragraph{Background derivatives of the fine kinetic form.}
The zero-background kinetic form is a curl on each fluctuation leg.
Its first background derivative is a sum of products with one finite
difference among three legs; its second background derivative needs none among
four legs. These representations have exponentially summable
coefficients with uniform bounds. For the first background derivative, subtract
its constant momentum coefficient, which is zero by the
Taylor calculation at scaling dimension four in Proposition~\ref{ch07:integrand-input}, and use
coordinate division. The zero-background representation is exactly
$d_1^*Hd_1$; the second background derivative already has the required form.
The number of small factors is
\[
 (n+2)+\text{number of differences}=4\qquad(0\le n\le2).
\]
H\"older with two fluctuation $L^2$ factors and the remaining supremum
factors gives $O(\mu^4)$ in normalized counting measure. Summing over the
fine base point contributes $L^4$, canceling it. A translation by $z$ introduces a
Bloch factor of modulus at most $C e^{\sigma|z|/L}$. The exponential
decay of the coefficients makes the resulting sum over $z$ uniformly bounded. Thus the fine kinetic Hessian and its first two
background derivatives have bounded operators between the stated
weighted spaces, uniformly in $L$. This is an estimate for the
left-coordinate response, whose definition by derivatives in left increments was used above.

\paragraph{The centered local gauge transformation and the penalty term.}
A small supremum bound alone is not enough for two fluctuation directions.
Use the uniform probability measure on the patch $Ly+\mathbb B_R$, $R=(j+3)L$, of the averaging operation and the three
norms $L^\infty,L^2,L^1$. Background directions use the first, one fluctuation
direction the second, and two fluctuation directions the third.
Products are controlled by $L^2L^2\subset L^1$, and averages contract
these norms. A contour of $O(L)$ links, each of whose positive derivatives
is $O(\mu)$ in its assigned norm, has bounded derivatives of every fixed
order. Indeed adjoining one more link changes the next derivative by
$O(\mu)$ times a polynomial in already bounded lower derivatives;
summing $O(L)$ such changes gives $O(1)$. For the coordinate contours of the transporters,
the link read at a fixed position sits at a translate of a point having some coordinates of each endpoint.
On the endpoint pair space this is a measure-preserving coordinate
permutation, followed by a bounded-measure restriction and translation.
The argument therefore applies to the averaged transporters.

Here is why centering gives uniform derivatives of the local gauge transformation. Let $F_x$ be the
minimizing local gauge transformation after recentering, $E_{xy}$ its averaged contour variable,
$\nu$ normalized counting measure on its block, and $\omega$ the
smoothing probability. If $\mathcal P$ is the nearest-group map and
$\Pi=D\mathcal P(1)$, define
\[
 J_0=\sum_y\nu(y)F_y,\qquad
 B_x=\sum_y\nu(y)F_y(E_{xy}-1),
\]
\[
 R(T)=\mathcal P(T)-\Pi T,\qquad
 S_x=\Pi B_x+R(J_0+B_x).
\]
The minimizing equation $F_x=\mathcal P(J_0+B_x)$ and centering
$\sum_x\omega(x)\Pi F_x=\Pi1$ imply exactly
\[
 F_x=\Pi1-\sum_y\omega(y)\Pi S_y+S_x.
\]
Now $E_{xy}(0)-1=0$ and $DR(1)=0$. At derivative order $n+1$ the
right side depends on derivatives of the local gauge transformation only through order $n$.
Induction using the preceding product bounds gives uniform derivatives
through order four, also with two fluctuation directions. The smoothing
probability satisfies $\omega\le C_j\nu$: its reduction modulo $L$
is uniform, so every atom has mass at most $L^{-4}$, and its support
has $O_j(L^4)$ sites. This controls the change of averaging measure.
The argument estimates the selected minimizer; its existence and local
uniqueness come from Chapter~\ref{ch04:blocking}.

The averaging of transporters and the fixed smooth penalty term are compositions of
these operations. Their mixed derivatives with up to two background and
two fluctuation directions are uniformly bounded. The
logarithmic derivative $X$ of the centered local gauge transformation has normalized $L^2$ output bounded by a
constant, so its unnormalized output is at most $CL^2$. Rooted and
centered descriptions differ by a triangular coordinate change of
determinant one; the same averaged-contour bound controls its background
coarse coordinates. Consequently $F=d_U^*P_f+m_\mu X$ has a uniform
weighted operator bound, since $m_\mu=O(\mu^2)$ cancels $L^2$.
Covariance of the local gauge transformation gives $X\mathcal R=I$ in the same coordinates, and
hence $F\mathcal R=\Delta_U+m_\mu$. Gauge signs are fixed here by
taking the fine component of $\mathcal R$ to be $d_U$.

\paragraph{Uniform inverse bounds.}
At zero background the regulated fine form has the weighted lower bound
$2^{-10}\min(1,\varepsilon)$, where $H_f\ge\varepsilon I$.
This follows from the exterior identity and the lower bound for the scalar regulator
proved in Proposition~\ref{ch07:integrand-input}. The kinetic form with the penalty
term added, the form $\mathcal G$ below, needs a quotient estimate. In centered coordinates it is
\[
 \mathcal G(u,r)=\langle u,K_fu\rangle
                         +b|M u-d_{0,c}r|^2.
\]
Here $M$ is the linear averaging operator. Indeed the rooted
linearization of the averaging operation is $Q=M+d_{0,c}R$ by \eqref{ch05:actual-Q}; setting
$r=r_{\rm root}-Ru$ converts $Qu-d_{0,c}r_{\rm root}$ to the
displayed expression. Under a fine gauge variation, the new scalar
coordinate changes by the smoothed scalar mean, since
$R d_{0,f}$ is root evaluation minus that mean. This is the same
centering used in the calculation for the local gauge transformation, and the triangular coordinate
change has determinant one.
Subtract a gauge vector to make every nonzero alias of $u$
transverse and eliminate $r$ using the scalar zero alias. The
latter scalar mean is $\mu^2$ times a product of centered box sums,
bounded away from zero after division by $\mu^2$ on the principal
real cube. Write the resulting fine aliases as $u_l$ and set
$A_h=\sum_{l\ne0}h_l^2|u_l|^2$. The lower bound on the fine metric gives control
of $A_h$, since $|q(k_l)|\ge c\mu$ for $l\ne0$.
The zero-alias coefficient of $M$ has least singular value $c_j\mu$:
its sum along the transporter is of order $L$ and the unitary block Fourier
coefficient is $L^{-2}$. For the upper bound on $M$, fix one of the $L$
steps of the transporter. Averaging the link at that step over the offsets gives coefficients
$O_j(L^{-4})$ on $O_j(L^4)$ sites, hence Euclidean row norm
$O_j(L^{-2})$. Summing over its $L$ positions along the transporter and converting
a weighted unit vector with $h_l\ge L^{-1}$ gives a fixed bound.
Thus the full $M$ has bounded weighted norm.
Consequently
\[
 \mu^2|u_0|^2\le C_j(|Mu|^2+A_h),
 \qquad
 \inf_a\|(u,r)-\mathcal R a\|_p^2\le C_j\mathcal G(u,r).
\]
These estimates include zero coarse momentum; the scalar zero
alias removes the coarse scalar coordinate there as well.

Moreover $\|\mathcal R a\|_p\le C|\Gamma_\mu a|$, where
$\Gamma_\mu=\Delta+m_\mu$. For the fine component this is
$h_l|q_l|\le h_l^2\le C\Gamma_\mu(k_l)$; the coarse averaging
component has norm $O(\mu^2)$ and obeys the same bound. Since
$F\mathcal R=\Gamma_\mu$ and $F$ is bounded, the quotient estimate
implies a positive weighted lower bound for $\mathcal G+F^*F$.
All constants are fixed before $L$.

The weights are not holomorphic. To continue this lower bound, at each real
base point hold $h_l$ constant while retaining the holomorphic block Fourier
matrix. On a fixed-radius coarse ball these constant weights are uniformly
comparable to the moving weights, because $h_l\ge\mu$ and
$|q(k_l(p))-q(k_l(x))|\le C\mu|p-x|$.
Cauchy's inequality in the coordinates with constant weights bounds the change in
operator norm, and a fixed smaller strip preserves half the real lower bound.
This proves uniform holomorphic inverse bounds for both matrices.
Differentiating $M^{-1}M=I$ then bounds the first two background derivatives of the inverses.

\paragraph{Trace norms and finite-volume coefficients.}
The kinetic response cancels identically, including its two derivatives,
in \eqref{ch07:theta}. The finite-rank terms retain their fixed-dimensional
coarse intermediate spaces. For the regulator, the argument given before this
subsection for the trace-class factor $P_\mu^3$ applies. In the terms containing
$m_\mu X$, its factor $\mu^2$ compensates for the factor $L^2$ of the derivative $X$ of the local gauge transformation.
In the separate term $P_f^*M_\mu P_f$, it compensates for two comparisons of weighted and Euclidean
norms, each of size $O(\mu^{-1})$. These large factors occur in
different terms; the same small factor is not used to compensate for both. Positive background derivatives of $d_U$ have a fine-link factor
$O(\mu)$; the single norm comparison contributes at most $O(\mu^{-1})$.
The norms of these first two background derivatives are therefore bounded in both
placements needed for $d_U^*d_U$ and $d_Ud_U^*$. The remaining
products use exactly these compatible input/output norms. Thus
$\Theta$ and its first two background derivatives have uniform trace norm. This proves all operator and trace bounds used in
Proposition~\ref{ch07:comparison-prop} for this same action and local gauge transformation.

Finally the infinite-lattice Fourier representations follow from the infinite
path coefficients and finite stencils of the averaging operation. All derivatives are
taken before restricting to a torus grid. Products and inverse derivatives
preserve this identification, with intermediate momenta shifted by the
inserted external momentum. The Fourier mesh estimate at the end of
Proposition~\ref{ch07:integrand-input} gives the finite-volume exponential
error. The preceding contact proof identifies the coefficient of
the Gaussian contribution \eqref{ch07:thermal}. Hence the comparison concerns the same density,
averaging operation and critical configuration throughout.


\begin{proposition}[Common comparison estimates]\label{ch07:comparison-input}
On a fixed positive strip of coarse momenta, the background derivatives at the critical
configuration and all intermediate shifted momenta obey the following
bounds, uniformly in $L$ and in the period. After removing the
contact term of the averaging operation, the blocked completed matrix $C_0$ and
$\widetilde A_\mu=\operatorname{diag}(A_\mu,I)$ have bounded inverses
and bounded first and second normalized background derivatives in
compatible weighted Hilbert norms. Their difference $\Theta$ and its
first two background derivatives have bounded trace norm in these same norms.
The infinite-lattice quadratic symbols represent the finite-torus coefficients
with an error tending exponentially to zero with the coarse period.
The linearization of the critical configuration is the one built with the averaging operation and obeys the
exact coarse--fine matching identities.
\end{proposition}

\begin{proof}
The preceding weighted inverse, background-derivative and trace estimates prove this
proposition for the normalized path action and the averaging operation.
\end{proof}

\begin{proposition}[Block comparison]\label{ch07:comparison-prop}
For the same normalized path action and averaging operation, the curvature-square
coefficients of the blocked response and the regulated fine determinant
satisfy $|\alpha_{N,m}-\alpha(\Pi_\mu)|\le C$, uniformly in $L,N,p$.
\end{proposition}

\begin{proof}
For a positive matrix continued holomorphically near zero background,
\[
 D_vD_w\log\det A
 =\tr(A^{-1}A_{vw}-A^{-1}A_vA^{-1}A_w).
\]
Subtract this formula for $C_0$ and $\widetilde A_\mu$, and use
$C_0^{-1}-\widetilde A_\mu^{-1}
=-C_0^{-1}\Theta\widetilde A_\mu^{-1}$.
Telescoping produces terms containing a trace-class difference or one
of its derivatives, with all other factors bounded. The trace ideal
inequality gives a uniform bound for the difference symbol on a common
strip. The ghosts cancel. The contact correction has zero $\alpha$ by
the preceding order argument.

Folding the fine symbol into coarse momentum also preserves $\alpha$.
The regulated determinant is gauge invariant, so its quadratic symbol
annihilates longitudinal legs. For a nonzero alias the equation
$K_f(k_l)N_l=A_lK_c(p)$ and the positive transverse fine inverse give\mgubi{Issue 10: $N_l$ is introduced only in the alias paragraph near the end of the chapter, and $A_l$ nowhere ($\mathcal A_l$ of Chapter~\ref{ch05:gaussian} is another object). Say that the alias amplitudes are those of $W_s$.}
$N_l^\perp=O(|p|^2)$, since $K_c(p)=O(|p|^2)$. Thus nonzero
aliases contribute order four even when their longitudinal Fourier coefficients
do not vanish. The zero alias has factors $L^4$ from
cell volume, $L^{-2}$ from the two interpolated link fields, and
$L^{-2}$ from extraction at momentum $p/L$. Their product is one.
The inherited loop has order four and contributes no $\alpha$.
Finally Cauchy's inequality on a circle inside the common strip bounds
the second Taylor coefficient of the difference by a fixed multiple of
its supremum. This proves the assertion.
\end{proof}

\subsection{The one-loop shift of the coupling}

\begin{proposition}[One-loop shift of the coupling]\label{ch07:slope}
For two consecutive scales with the normalized path action, the lower bound on the metric and the
identity \eqref{ch05:stationary-tangent} used in Section~\ref{ch07:weighted},
\[
 |\alpha_{N,m}+\kappa_G\log L|\le C
\]
with a constant depending only on the fixed stage inputs.
\end{proposition}

\begin{proof}
On the unit three-sphere the angular means are
$\mathbb E\theta_2^2=1/4$ and
$\mathbb E\theta_1^2\theta_2^2=1/24$. Therefore
\[
 (2\pi)^{-4}\int_{\mu<|k|<1}h(k)\,dk
 =\frac{8\pi^2C_G}{(2\pi)^4}
   \left(\frac14-\frac16-1\right)\log\mu^{-1}
 =-\kappa_G\log\mu^{-1}.
\]
The two error integrals on this annulus are uniformly bounded:
$\int_\mu^1 r^3r^{-3}dr\le1$ and
$\mu^2\int_\mu^1r^3r^{-6}dr\le1/2$.
The ball $|k|\le\mu$ has volume $O(\mu^4)$ and integrand
$O(\mu^{-4})$; the region outside the unit ball is also uniformly
integrable. Thus $\alpha(\Pi_\mu)=-\kappa_G\log\mu^{-1}+O(1)$.
Proposition~\ref{ch07:comparison-prop} and $\mu=L^{-1}$ finish the proof.
\end{proof}

The sign is for the step from scale $m+1$ to scale $m$: the coefficient $\beta$ of the
kinetic action decreases by $\kappa_G\log L$ at one loop.
This proposition does not assert the exact nonlinear recursion for
$\beta$. Connected interactions contribute to that recursion and are
handled by the one-step integral identity.

\section{Classical matching through fourth order}\label{ch07:classical}

\subsection{Uniform classical coefficients at the critical configuration}
\label{ch07:classical-uniform}

To construct the cubic and quartic matching terms, we first need
classical coefficient bounds fixed before $L$. The fixed-$L$ chart for the critical configuration of
Chapter~\ref{ch05:gaussian} provides existence and differentiation, but
its inverse in the block axial gauge alone does not provide these stronger bounds.
We obtain them from the weighted quotient just constructed.

Let $\mathscr S(v,z)$ be the joint classical action, with coarse field
$v$ and fluctuation (internal) coordinates $z$, the fine field together with the local gauge transformation. It contains the fine
kinetic action and the penalty term built with the averaging operation, without $\beta$ or the
inherited loop. At zero its first derivative vanishes. Write
$\mathscr H=D_z^2\mathscr S(0)$ and let $\mathcal R$ denote its
infinitesimal gauge action on the fluctuation coordinates. Its gauge-fixed inverse is
\[
 \mathscr C^{-1}=(\mathscr H+F^*F)^{-1},
\]
with the uniform weighted bound just proved. For a coarse direction $u$
let $T u=(u,W_a u,r_u)$ be the full linearization of the critical configuration in the smooth
centered representative. In particular
\begin{equation}\label{ch07:tangent-orthogonal}
 D^2\mathscr S(0)[Tu,(0,z)]=0
 \quad\text{for every internal }z.
\end{equation}
This is the differentiated stationarity equation, with the
identification for the averaging operation in Chapter~\ref{ch05:gaussian}. Stationarity
in all joint fluctuation directions follows from stationarity on the
slice of the block axial gauge and gauge invariance: that slice together with the joint
gauge directions spans the tangent space of the fluctuation coordinates. Thus coarse-frame
directions are included in \eqref{ch07:tangent-orthogonal}.
Define the covector on fluctuation directions
\[
 B(u,v)[z]=D^3\mathscr S(0)[Tu,Tv,(0,z)].
\]

Both the cubic and quartic vertices with background legs and $B$ have uniformly
bounded infinite-lattice Fourier kernels. Here is the power count in their
weighted norms. A fine-action cubic vertex has at least one
finite difference among its three legs; a quartic vertex needs none.
Their exponentially summable coefficient bounds follow from the same
expansion of the normalized path action and coordinate division used for the background
derivatives of the kinetic form above. An undifferentiated background leg contributes $O(\mu)$ in
supremum norm, and its differentiated leg contributes $O(\mu^2)$. A fluctuation
leg of unit weighted norm contributes these same powers in normalized $L^2$ norm.
Thus three background legs with one difference, four background legs,
and two background legs with one fluctuation leg and one difference each
contribute $O(\mu^4)$ in normalized summation over positions within the block. In the last case
$L^2\subset L^1$ on probability measure is sufficient. The sum over
$L^4$ fine base points cancels this factor. Finitely many coordinate-change
terms have the same or higher scaling dimension. Penalty vertices are bounded by
the induction on mixed derivatives of the centered local gauge transformation above, which allows up to
four derivatives and includes the bounded coarse scalar legs $r_u$.
The bounds for $B$ are dual weighted norms, not Euclidean coordinate
bounds with a dimension-dependent norm comparison.

The covector $B$ annihilates gauge directions. To check this, write the
gauge identity in the fluctuation coordinates as
$D_z\mathscr S(v,z)[\mathcal R(v,z)a]=0$ and differentiate twice along
$Tu,Tv$. The term containing $D\mathscr S(0)$ is zero. The two terms
with a derivative of $\mathcal R$ vanish by
\eqref{ch07:tangent-orthogonal}, since those derivatives are fluctuation
vectors. The remaining term is
\[
                       B(u,v)[\mathcal R(0)a]=0.
\]
Therefore $w=\mathscr C^{-1}B(u,v)$ is a representative of the quotient
solution $\mathscr H w=B(u,v)$. Indeed, pairing
$\mathscr Cw=B(u,v)$ with $\mathcal R a$ gives
$\langle F\mathcal R a,Fw\rangle=0$. The regulated gauge matrix
$F\mathcal R$ is invertible, so $Fw=0$ and $\mathscr H w=B(u,v)$.
Any other solution differs by a gauge vector, against which $B$ vanishes.
This justifies using the bounded gauge-fixed inverse in the exchange term;
it does not replace a degenerate Hessian by an arbitrary inverse.

If $E(v)$ is the value at the critical configuration, implicit differentiation now gives
\begin{align}\label{ch07:uniform-stationary-jets}
 D^3E(0)[u_1,u_2,u_3]
   &=D^3\mathscr S(0)[Tu_1,Tu_2,Tu_3],\\
 D^4E(0)[u_1,u_2,u_3,u_4]
   &=D^4\mathscr S(0)[Tu_1,Tu_2,Tu_3,Tu_4]\notag\\
   &\quad-\sum_{\{ij\}\{kl\}}
          B(u_i,u_j)\mathscr C^{-1}B(u_k,u_l),\notag
\end{align}
where the sum is over the three partitions of $\{1,2,3,4\}$ into pairs.
Each summand evaluates the covector $B(u_i,u_j)$ on the vector
$\mathscr C^{-1}B(u_k,u_l)$, using the counting inner product to
identify the latter covector with a vector before applying the inverse.
These formulas follow by differentiating stationarity; the higher
derivatives of the critical-configuration map multiply its first derivative or
\eqref{ch07:tangent-orthogonal}. On the diagonal the last sum is
$3B(u,u)\mathscr C^{-1}B(u,u)$, agreeing with the earlier formula.
The gauge-annihilation argument makes the formulas independent of the
representative chosen to differentiate the critical-configuration map.

For Fourier waves, each exchange inverse has intermediate coarse
momentum equal to a sum of two external momenta. Restrict each external
strip by a fixed factor so that all these sums remain in the inverse
strip. The vertex and inverse bounds then give a common strip and a
uniform bound for both symbols in
\eqref{ch07:uniform-stationary-jets}, before choosing $L$.
Their finite-volume identities are exact: every vertex and inverse is
the restriction of the same infinite-lattice Fourier matrix, and momentum
conservation leaves the coarse volume factor exactly once. Dividing by
$3!$ and $4!$ gives normalized infinite-lattice Taylor symbols with the required
bounds. No neighborhood for the critical configuration of
size uniform in $L$ has been inferred from this argument; only its
fixed-order coefficients need that uniformity here.

\subsection{Construction of the matching curvature terms}

Let $E(v)$ be the classical action at the critical configuration of Chapter~5, with the
factor $\beta$ removed. Its Hessian is prescribed by the exact free
block relation. Write the corresponding plaquette metric as
$H=c_{\rm dir}I+\rho^*\rho$, with $H(0)=I$ and $c_{\rm dir}>0$.
The direct Wilson term together with the squares of the transported
rows of $\rho$ gives an analytic gauge-invariant action $J(v)$ with
this Hessian. Indeed the transported rows linearize to $\rho d_1v$,
so
\[
 D^2J(0)[v,v]=\langle d_1v,Hd_1v\rangle=D^2E(0)[v,v].
\]
The construction uses the coarse metric matched to the incoming
Gaussian form, not an arbitrary positive coarse metric.

The preceding coefficient estimates at the critical configuration and the fixed row bounds
give the hypotheses of the following construction.

\begin{proposition}[Classical matching]
\label{ch07:classical-match}
Suppose $E$ and $J$ are invariant analytic germs on every torus, have
zero value and first derivative and identical Hessians. Suppose their
normalized coefficients of orders three and four have infinite-lattice symbols with strip
bounds, and the row coefficients of $J$ have an exponential bound.
There are invariant cubic and quartic curvature terms $S_3,S_4$, chosen
once for all periods, such that
\[
 E-J-S_3-S_4=O(v^5).
\]
Their weighted sums are bounded by a fixed multiple of the input Taylor
bounds on a smaller positive width.
\end{proposition}

\begin{proof}
The leading coefficient of $E-J$ is cubic and satisfies the homogeneous
Ward identities. Lemma~\ref{ch07:curl} and the transported construction
produce $S_3$ with that coefficient. The fourth coefficient of
$E-J-S_3$ is again transverse, since its lower coefficients vanish.
Apply the lemma again to obtain $S_4$. It is essential to subtract the
fourth coefficient of $S_3$: a nonlinear transported cubic curvature
term has nonzero higher Taylor coefficients. The exponential path
bound controls each fixed Taylor derivative on a smaller complex link
ball; factors of path length are absorbed by a width loss. Thus the
fourth coefficient of $S_3$ has the required common bound. This proves
both the matching and the bounds. Polarization converts the equality on
rays to equality of all derivatives through order four.
\end{proof}

\subsection{Preserving the classical coefficient bound}

These estimates can also be arranged to preserve a fixed classical
coefficient bound, rather than producing an unspecified function of the
incoming bound. Keep the Gaussian metric bounds fixed, and denote the
bound on the remaining incoming path coefficients by $A$. Compare the
fine action with the Wilson action, evaluated on the \emph{same}
linearizations $T$ and with the \emph{same} gauge-fixed inverse $\mathscr C^{-1}$.
The normalization at scaling dimension four proved above implies that their cubic
vertex difference has at least two finite differences, and their
quartic vertex difference has at least one. Coordinate division of their
infinite-lattice symbols gives these factors with exponentially summable
coefficients bounded by a polynomial $P(A)$. A slightly stronger incoming
exponential width absorbs the polynomial path-length factors.

When all legs are background legs, the extra difference gives $L^4\mu^5=\mu$.
For the difference of the covectors $B$, move differences on the single
fluctuation leg to the two background legs by momentum conservation. In a
coordinate this uses
\[
 q(-p_1-p_2)=-e^{-i(p_1+p_2)}
                \bigl(q(p_1)+e^{ip_1}q(p_2)\bigr).
\]
It is a discrete integration by parts, with bounded translation factors;
it leaves at most two differences on any background leg. The
bounds on the smooth representative $W_a$ through order two therefore give the same gain.
The fluctuation leg is estimated only in normalized $L^2$, without asking
for a second weighted derivative. Thus the vertex and dual-vertex
remainders are bounded by $P(A)/L$.

The Wilson vertices, penalty term, and their exchange pairing have
bounds depending only on the fixed Gaussian metric and the data of the averaging operation.
The inverse in the exchange term depends only on the quadratic action,
so higher incoming coefficients do not change it. Expanding the two
factors $B=B_{\rm W}+\delta B$ in
\eqref{ch07:uniform-stationary-jets} shows that the normalized classical
symbols, and then their curvature families, satisfy a bound of the form
\begin{equation}\label{ch07:classical-budget}
             A_{\rm out}\le a_{\rm stage}+\frac{P_{\rm stage}(A)}{L}.
\end{equation}
The squared remainder contributes $P(A)^2/L^2$ and is absorbed into the displayed
polynomial. Subtracting the prescribed coarse quadratic row and the
quartic contribution of the constructed cubic family preserves this
form: those operations have fixed coarse-scale bounds, and the latter
operation is linear in that cubic family.
The output strip is chosen from the Gaussian coarse strip. It need not
shrink with each step: an incoming coefficient displaced by $z$ contributes a factor
at most $e^{c\sigma|z|/L}$ on that output strip, which is absorbed by the
fixed incoming width after choosing $L$ large enough. Finally choose an
outgoing classical bound larger than $a_{\rm stage}$ and then $L$
so that the second term in \eqref{ch07:classical-budget} is smaller than the difference between these two bounds.
For a finite collection of stage data these choices are simultaneous.

Consequently the cubic and quartic curvature families are constructed
for this same action at the critical configuration with bounds fixed at the stage before
$L$. Their input Hessian is the matched coarse metric.
This completes the classical matching construction without invoking the schedule of Chapter~\ref{ch10:trajectory}
or a nonlinear local-integral estimate.

\section{Inherited loops: four differences and alias suppression}

Taylor degree here counts field factors. It is distinct from scaling dimension, which also counts derivatives, as in Section~\ref{ch02:jets-section}; a momentum difference therefore improves scaling dimension without adding a field factor.
\label{ch07:inherited-proof}

Let $\Lambda_{m+1}^{\rm loop}$ denote the incoming weighted loop sum.
The estimates of Chapter~\ref{ch05:gaussian} for the representatives of the linearization of the critical configuration give, with ``pulled back'' meaning composed with the critical-configuration map,
\begin{equation}\label{ch07:inherited}
 \|\mathcal L_{m+1}\hbox{ pulled back}\|_{w'}
 \le C L^{-2}\Lambda_{m+1}^{\rm loop}
\end{equation}
To prove this bound, group incoming quadratic curvature terms by their relative displacement
$z$. Their coefficient matrix $r(z)$ has exponential sum bounded by
$\Lambda_{m+1}^{\rm loop}$ and zero zeroth and first moments. Thus the
fine link symbol
\[
 S(p)=d_1(-p)^{\mathsf T}\widehat r(p)d_1(p)
\]
has order four, satisfies $S d_0=d_0^*S=0$, and has strip bound
$C\Lambda_{m+1}^{\rm loop}$. This bound holds componentwise in the
Lie algebra and hence also for the symmetrized bilinear form, with a
constant depending on $\dim\gLie$.

An analytic periodic function with an order-four zero is a sum of
products of four coordinate differences $q_a(p)$ times analytic periodic
functions on a smaller strip, with linearly bounded coefficient norms.
To see this quantitatively, write
$f(p)-f(0)=\sum_a q_a(p)D_a f(p_1,\ldots,p_a,0,\ldots,0)$
using the divisions of Lemma~\ref{ch07:curl}, and repeat on the quotients,
which retain one fewer vanishing order. There are finitely many terms.
For a quadratic symbol, move two differences to the opposite leg by
$q_a(p)=-e^{ip_a}q_a(-p)$. The phases merely translate a coefficient
by finitely many lattice sites. Each term now has two differences on
each of its two field arguments.

The representative $W_s$, built with the linear averaging operator, has exponentially local infinite-lattice
kernels with $\nabla_f^rW_s=O(L^{-1-r})$ for $r\le2$, by
Section~\ref{ch05:transports}. By \eqref{ch05:corrected-transport}, the linearization built with the averaging operation differs
from this representative by a fine gradient; the gauge-invariant
quadratic symbol $S$ annihilates that difference on both sides.
Consequently it is legitimate to use $W_s$ in this calculation, even
if the representative $W_F$ in the block axial gauge itself has poor difference bounds.
Each differentiated leg contributes $CL^{-3}$. Summing over the fine base point in
one coarse block contributes $L^4$, so the pulled-back quadratic symbol has
bound
\[
 C L^4 L^{-3}L^{-3}\Lambda_{m+1}^{\rm loop}
       =CL^{-2}\Lambda_{m+1}^{\rm loop}.
\]
Displacing an argument by a fine displacement $z$ contributes a factor at most
$Ce^{\sigma|z|/L}$ on a coarse strip of width $\sigma$.
Choose $L$ after the fine coefficient width $w$ so that $\sigma/L<w$;
absolute summability controls these displacements. Holomorphy and the
infinite-lattice representation of the finite-torus coefficients follow from normal convergence.
A final fixed strip reduction converts this to the weighted coefficient
bound required in \eqref{ch07:inherited}.

The order-four zero on the coarse lattice requires another argument.
For aliases $k_l(p)=(p+2\pi l)/L$, let $N_l(p)$ be the Fourier
coefficient at the alias $l$ of the linearization of the critical configuration. The coarse symbol is
\[
 H(p)=L^4\sum_l N_{-l}(-p)^{\mathsf T}S(k_l(p))N_l(p).
\]
For $l\ne0$, use $K_f(k_l)N_l=A_lK_c(p)$, the exact equation for
that linearization. Project $N_l$ orthogonally off the gradient
$q(k_l)$. The positive fine curvature metric implies that
$K_f(k_l)+d_0(k_l)d_0(k_l)^*$ is invertible near $p=0$ for this fixed
$l,L$. As $K_c(p)=O(|p|^2)$, the transverse coefficient is $O(|p|^2)$.
The Ward identities of $S$ remove the longitudinal coefficients, so each
nonzero alias has order four. The zero alias has order four
because $S(p/L)$ does. This argument requires no lower inverse bound
uniform in $L$ for an individual nonzero alias.

Finally the pullback by the nonlinear critical-configuration map has the same Hessian as this
linear pullback: the inherited action has zero first derivative, so its
chain rule has no term involving the second derivative of the critical-configuration map.
This proves the weighted gain and retained moment cancellations for the
inherited action, once the linearization built with the averaging operation and the representative $W_s$
have been identified up to gauge as above.

\section{Output coefficients and the remaining analytic width}\label{ch07:output}

Suppose the determinant symbol at the critical configuration has bound $\Phi$ at a common
width, the inherited loop symbol has bound $R$ and order four, and the
kinetic symbol has bound $B_H$. Subtract the extracted coefficient
$\alpha b_H$. Lemma~\ref{ch07:unique} and the two-leg division produce
a quadratic curvature family with bound
\[
 K(\Phi+|\alpha|B_H+R)
\]
and the zero moments \eqref{ch07:moments}. Together with the prescribed
rows and the cubic/quartic families of
Proposition~\ref{ch07:classical-match}, this defines one common output
path action. If the infinite-lattice determinant symbol approximates its finite
period values with error $|\Lambda|\epsilon$, the same coefficient
error is retained by linear subtraction. For real plane waves with
component amplitudes at most one, the Hessian error is at most
$2(4\dim\gLie)^2|\Lambda|\epsilon$: expansion in component waves
contributes the square of their number, and $D^2=2C_2$.

All sums should be bounded first at a width $w'>w>0$.
Each fixed differentiation of a transported term contributes a polynomial factor in
its path length; the factor $e^{-(w'-w)\operatorname{length}}$ absorbs
it. This gives the fixed derivative orders needed for local expansion.
Parameter derivatives require the corresponding derivative bounds on
the infinite-lattice input kernels. Linear curvature division preserves such
bounds because its operations and widths are chosen independently of
the particular kernel. No parameter regularity follows merely from
matching the kernels pointwise.

For a quantitative descent, use the inherited estimate
\eqref{ch07:inherited}. If the new loop
bound is $A L^4+B\log L+C_0$, choose $L$ so that $CL^{-2}\le1/4$,
then the lower bound on the reference coupling so that
\[
 A L^4+B\log L+C_0\le\tfrac12\Lambda_m^{\rm loop},\qquad
 \Lambda_{m+1}^{\rm loop}\le2\Lambda_m^{\rm loop}.
\]
Their sum is at most the outgoing bound. This is the order of choices:
widths and fixed amplitudes, block ratio, then the lower bound on the reference coupling.

The comparison bounds and lattice integrand estimates have now been
derived from the same normalized analytic path action, positive metric and
averaging operation at the critical configuration. The inherited estimate follows from the
identified linearization of the critical configuration and the bounds on the infinite-lattice kernels of $W_s$. These results are
one-step statements: their application to an entire selected history
requires that its incoming path actions obey the fixed stage bounds and its
charts for the critical configuration satisfy the identifications of Chapter~\ref{ch05:gaussian}.
They do not replace the nonlinear local-integral bounds or prove the
existence of an infinite sequence of scales.
